\documentclass[../main.tex]{subfiles}

\begin{document}

\subsection{Notation}

$\N$ denotes the set of natural numbers, $\N_0 := \N \cup \{0\}$, and $\N_2 := \N \backslash \{1\}$. $\Q$ stands for the set of rational numbers, $\R$ for the set of real numbers and $\R_+$ for the set of nonnegative real numbers. For a set $A$, $\card A$ denotes its cardinality. If $\card X = \infty$, we may write $\card X = \aleph_0$ if $X$ is countable, and $\card X = 2^{\aleph_0}$ if $X$ is uncountable. For every subset $A$ of $\R$, $\mathring A$ denotes the interior of $A$. Specifically, if $A = [a,b]$, then $\mathring A = (a,b)$, for all $a,b \in \R$, $a<b$. Let $f : X \to X$ be a map from a set $X$ into itself. Then, $f^0$ denotes the identity map on $X$, $f^1 := f$, and $f^{n+1} := f \circ f^n$, for $n \in \N$. Let $f : X \to Y$ be a map from a set $X$ into a set $Y$. For a given subset $A \subset Y$, we let $f(A) := \{ f(x) : x \in A\}$ and $f^{-1}(A) := \{ x \in X : f(x) \in A\}$.

Let $A$ be a set. $A^\omega$ denotes the set of infinite sequences with elements in $A$, $A^\ast$ denotes the set of sequences of finite, but otherwise arbitrary length, of elements in $A$, and $A^n$ denotes the set of sequences of elements in $A$ of length $n \in \N$. We adopt the notations $\xf = \xf_1 \xf_2 \ldots \in A^\omega$, and $x = x_1 \ldots \vspace{0.2mm} x_n \in A^\ast$, for $n \in \N$. $\epsilon$ denotes the empty sequence. For a sequence $x \in A^\ast$, $x_n$ denotes the $n$-th element of $x$, $x_n$ is defined to be $x_{1} \ldots x_n$, and $x|n^m$ stands for the sequence $x_{n} x_{n+1} \ldots x_m$. Similarly, for $\xf \in A^\omega$, $\xf_n$ denotes the $n$-th element of the sequence $\xf$, $\xf|_n$ is defined to be $\xf_{1} \ldots \xf\pn n$, and $\xf\pnm{n}{m}$ stands for the finite sequence $\xf_{n} \xf_{n+1} \ldots \xf_m$. $\len x$ refers to the length of $x$, i.e. for $x \in A^n$, we have $\len x = n$, and by convention, we set $\len \xf = \infty$ for $\xf \in A^\omega$. For $x,y \in A^\ast$, $x y$ denotes the concatenation of $x$ and $y$, and likewise for $x \yf$, with $x \in A^\ast$ and $\yf \in A^\omega$. For $x^{(1)},\ldots, x^{(n)} \in A^\ast$, $\coprod_{i=1}^n x^{(i)}$ stands for the concatenation $x^{(1)} x^{(2)} \ldots \hspace{0.2mm} x^{(n)}$. For an infinite family $(x^{(i)})_{i\in \N} \in A^\ast$, $\coprod_{i=1}^\infty x^{(i)} \in A^\omega$ denotes the concatenation of the elements in $(x^{(i)})_{i \in \N}$. Finally, for $x \in A^\ast$, we set $x^0 = \epsilon$, $x^n := \coprod_{i=1}^n x$, for $n \in \N$, and $x^\infty := \coprod_{i=1}^\infty x$, to denote the $n$-fold and infinite repetitions, respectively. For instance, $0^\infty$ denotes the infinite sequence containing only zeroes, and $1^5 = 11111$. We may combine repetitions and concatenations in a single expression, as for example $xy^5 = xyyyyy$, for $x,y \in \{0,1\}^\ast$. So that there is no confusion, we may use parentheses. For instance, $(10)^5 = 1010101010$, while $10^5 = 100000$. $\{0,1\}^\ast$ is the set of finite binary sequences, and $\{0,1\}^\omega$ the set of infinite binary sequences. Following \cite[Ch. 1.4]{li2008introduction}, we map $\{0,1\}^\ast$ one-to-one onto $\N_0$ by indexing each string through the quasi-lexicographical ordering \cite[Section 1.1]{calude2002information}, i.e.,
\begin{equation}\label{eq:quasilexicographic enumeration}
	(\epsilon,0), (0,1), (1,2), (00,3), (01,4), (10,5), (11,6), \ldots,
\end{equation}
where each of the above pairs contains first an element $x \in \{0,1\}^\ast$ and second the index of $x$ in the quasi-lexicographical ordering of $\{0,1\}^\ast$. We make the confusion between $x \in \{0,1\}^\ast$ and its index, so that we can use them interchangeably in equations and formulae. For $x \in \{0,1\}^\ast$, we define $\bar x := 1^{|x|}0x$, which is often referred to as prefixing $x$, i.e. $\{\bar x : x \in \{0,1\}^\ast\}$ forms a prefix-free set. For $x,y \in \{0,1\}^\ast$, we set $\langle x,y \rangle := \bar x y$. Note that $|\langle x,y \rangle| = 2 |x| + |y| + 1$. For $q \in \N_2$, $A_q := \{0, \ldots, q-1\}$, $A_q^\ast$ refers to the set of finite $q$-ary sequences, $A_q^\omega$ is the set of infinite $q$-ary sequences. For $\xf,\yf \in \{0, \ldots, q-1\}^\omega$, we define the lexicographical ordering as $\xf <_L \yf$ to express that there exists $j \in \N$ such that $\xf_i = \yf_i$ for all $i < j$, and $\xf_j < \yf_j$.


\bibliographystyle{IEEEtranS}
\bibliography{biblio.bib}

\end{document}